\chapter*{Organización del trabajo en grupo}
\addcontentsline{toc}{chapter}{Organización del trabajo en grupo}
\setcounter{secnumdepth}{-1}

\section{Introducción al trabajo grupal}

El presente Trabajo Fin de Máster (TFM) se ha desarrollado en modalidad grupal, de conformidad con el procedimiento establecido por la Universidad Internacional de La Rioja (UNIR). Esta modalidad resulta coherente con la naturaleza del proyecto: se diseña, implementa y evalúa un sistema inteligente de apoyo a la decisión para la priorización temprana de incidentes 112 en Castilla y León, combinando procesamiento del lenguaje natural, aprendizaje interpretable, explicación normativa y desarrollo de software.

La organización del equipo se articula en tres capas funcionales y una función transversal de integración. La Capa 1 estructura el texto operativo en variables y señales predecisionales. La Capa 2, desarrollada como bloque metodológico central, transforma esas señales en una recomendación P1--P4 interpretable y evaluable. La Capa 3 convierte la recomendación en una explicación operativa mediante un modelo de lenguaje de gran tamaño (\textit{Large Language Model}, LLM) local, generación aumentada por recuperación (\textit{Retrieval-Augmented Generation}, RAG) y herramientas basadas en \textit{Model Context Protocol} (MCP).

El reparto no se limita a distribuir tareas de programación. Cada integrante asume una contribución diferenciada al componente inteligente del sistema, de modo que el resultado conjunto pueda defenderse como una solución integrada y no como una suma de piezas independientes.

\section{Integrantes del equipo y roles asignados}

El equipo está formado por tres estudiantes del Máster Universitario en Inteligencia Artificial de la UNIR, bajo la dirección del Dr. Andrés Soto Villaverde. Cada integrante aporta un perfil complementario y una responsabilidad principal dentro de la arquitectura.

\subsection{Capa 1: extracción de variables operativas}
\textbf{Responsable principal:} Ancor González Carballo.

Su bloque de trabajo se orienta a la transformación del texto libre del incidente en una representación estructurada compatible con la priorización posterior. Esta capa combina procesamiento del lenguaje natural auditable, detección de señales operativas, tratamiento de negaciones y normalización de variables V01--V15. En la versión evaluada se presenta como un extractor determinista y reproducible; la incorporación de clasificadores aprendidos queda como evolución futura sujeta a una validación específica. Esta contribución es importante porque la calidad de la priorización depende de que el sistema represente correctamente el incidente antes de aplicar reglas o modelos.

\subsection{Capa 2: supervisión débil, RuleFit y evaluación}
\textbf{Responsable principal:} Juan Carlos Albert Fillol.

Su bloque de trabajo se centra en convertir datos reales e incompletos del Registro 112 Castilla y León en una recomendación de prioridad interpretable, evaluable y trazable. Para ello formaliza la escala académica P1--P4, audita el conjunto de datos, construye etiquetas mediante supervisión débil, define una línea base heurística, selecciona RuleFit-lite como motor interpretable y evalúa el comportamiento del sistema mediante macro-F1, sensibilidad de P1, estabilidad temporal, análisis por subgrupos y errores críticos.

Juan Carlos asume, además, la coordinación metodológica y el cierre del bloque de datos, supervisión débil, línea base heurística, priorización interpretable, evaluación, revisión interna y trazabilidad de Capa 2. Su aportación conecta la salida de Capa 1 con la explicación de Capa 3, sin asumir la autoría de las capas desarrolladas por los demás integrantes.

\subsection{Capa 3: explicación generativa y recuperación normativa}
\textbf{Responsable principal:} Brian Mathias Pesci Juliani.

Su bloque de trabajo comprende la explicación y la integración conceptual del prototipo. Incluye la recuperación normativa, la explicación generativa local y la trazabilidad de la inferencia. Su contribución consiste en convertir la salida del modelo interpretable en una recomendación comprensible y revisable, sin modificar la prioridad calculada por Capa 2 y manteniendo la ejecución local de los datos sensibles.

\section{Distribución de capítulos y anexos}

La Tabla~\ref{tab:tab1secCTF} recoge la distribución de responsabilidades sobre la memoria. El reparto final se ha definido con un único responsable por archivo para evitar duplicidades, solapamientos y conflictos de integración. De esta forma, cada integrante revisa su parte, identifica los cambios necesarios y aporta texto, evidencias, capturas o referencias para la unificación final.

\begin{table}[htbp]
\centering
\scriptsize
\setlength{\tabcolsep}{4pt}
\begin{tabular}{|p{4.2cm}|p{4.2cm}|p{4.6cm}|}
\hline
\textbf{Apartado / capítulo} & \textbf{Responsables} & \textbf{Tipo de contribución} \\ \hline
\texttt{chap0.tex} & Juan Carlos & Organización del trabajo y reparto final \\ \hline
\texttt{chap1.tex} & Juan Carlos & Introducción, problema, alcance y contribución IA \\ \hline
\texttt{chap2.tex} & Brian & Contexto, estado del arte y bibliografía APA 7 \\ \hline
\texttt{chap3.tex} & Ancor & Marco operativo, normativo y contexto 112 CyL \\ \hline
\texttt{chap4.tex} & Juan Carlos & Objetivos, preguntas de investigación y metodología \\ \hline
\texttt{chap5.tex} & Ancor & Requisitos y especificación funcional \\ \hline
\texttt{chap6.tex} & Brian & Arquitectura, contratos y diseño del sistema \\ \hline
\texttt{chap7.tex} & Juan Carlos & Datos de Castilla y León, supervisión débil, prevención de fugas y particiones \\ \hline
\texttt{chap8.tex} & Brian & Integración del prototipo y demostración funcional \\ \hline
\texttt{chap9.tex} & Juan Carlos & Evaluación de Capa 2, integración, revisión preparada y análisis por subgrupos \\ \hline
\texttt{chap10.tex} & Ancor & Conclusiones, limitaciones y trabajo futuro \\ \hline
\texttt{anexo\_a.tex}, \texttt{anexo\_b.tex}, \texttt{anexo\_g.tex} & Ancor & Taxonomía, variables y guía de uso \\ \hline
\texttt{anexo\_c.tex}, \texttt{anexo\_d.tex}, \texttt{anexo\_e.tex} & Juan Carlos & Guía P1--P4, protocolo de revisión y línea base heurística \\ \hline
\texttt{anexo\_f.tex}, \texttt{anexo\_h.tex}, \texttt{anexo\_l.tex}; \texttt{bibliografia.bib} & Brian & Correspondencia conceptual-software, transparencia, corpus normativo consolidado, fichas de modelo y bibliografía \\ \hline
\end{tabular}
\caption{Distribución de capítulos, responsables y tipo de contribución.}
\label{tab:tab1secCTF}
\smallskip
{\small\textit{Fuente: elaboración propia.}\par}
\end{table}

\section{Metodología de coordinación del grupo}

La coordinación se ha apoyado en reuniones periódicas, control de versiones y una planificación por tareas trazables. Las decisiones de alcance se documentan de forma común y los contratos de datos permiten que cada capa avance sin romper la integración.

El equipo ha trabajado con una lógica iterativa e incremental. Primero se fijaron los acuerdos de integración y la constitución del proyecto; después se definieron el conjunto de datos, la guía de etiquetado y el corpus normativo; posteriormente se implementaron las capas técnica y explicativa; y finalmente se orienta el trabajo hacia la evaluación, la redacción y la defensa. Esta organización permite mantener alineados el prototipo software, la memoria académica y la evidencia reproducible.

De este modo, la organización del trabajo refleja la idea central del TFM: diseñar, implementar y evaluar un sistema de apoyo a la decisión para incidentes 112 en Castilla y León, con prioridad explicable P1--P4, supervisión humana, soberanía del dato y trazabilidad normativa.

\setcounter{secnumdepth}{2}
